\usetheme{Madrid}
\usecolortheme{seahorse}
\setbeamertemplate{navigation symbols}{}
\setbeamertemplate{itemize items}[circle]
\setbeamertemplate{enumerate items}[default]

\usepackage{amsmath,amssymb,mathtools}

\newcommand{\R}{\mathbb{R}}
\newcommand{\N}{\mathbb{N}}
\newcommand{\Prob}{\mathbb{P}}
\newcommand{\E}{\mathbb{E}}
\newcommand{\var}{\operatorname{var}}
\newcommand{\one}{\mathbf{1}}
\newcommand{\eps}{\varepsilon}
\newcommand{\norm}[1]{\left\lVert #1\right\rVert}
\newcommand{\vikiname}[1]{\par\smallskip{\scriptsize\color{blue!60!black}\textbf{LLMwiki:} \texttt{\detokenize{#1}}}}
\newcommand{\blankproof}{%
  \vspace{2mm}
  {\color{gray!55}\rule{\linewidth}{0.4pt}}\par\vspace{5mm}
  {\color{gray!55}\rule{\linewidth}{0.4pt}}\par\vspace{5mm}
  {\color{gray!55}\rule{\linewidth}{0.4pt}}\par\vspace{5mm}
  {\color{gray!55}\rule{\linewidth}{0.4pt}}\par
}
\newcommand{\longblankproof}{%
  \vspace{1mm}
  {\color{gray!55}\rule{\linewidth}{0.4pt}}\par\vspace{4mm}
  {\color{gray!55}\rule{\linewidth}{0.4pt}}\par\vspace{4mm}
  {\color{gray!55}\rule{\linewidth}{0.4pt}}\par\vspace{4mm}
  {\color{gray!55}\rule{\linewidth}{0.4pt}}\par\vspace{4mm}
  {\color{gray!55}\rule{\linewidth}{0.4pt}}\par\vspace{4mm}
  {\color{gray!55}\rule{\linewidth}{0.4pt}}\par
}
\newcommand{\proofblock}[1]{\ifteacher #1 \else \blankproof \fi}
\newcommand{\longproofblock}[1]{\ifteacher #1 \else \longblankproof \fi}

\title[Chapter 1.6]{Chapter 1.6 Expected Value}
\subtitle{Rick Durrett Probability: Theory and Examples}
\author{Hu Jingwu}
\date{}

\begin{document}

\begin{frame}
  \titlepage
  \vspace{-5mm}
  \begin{center}
    \small Built from the local Chapter 1 LLMwiki flash cards.
  \end{center}
\end{frame}

\begin{frame}{Roadmap for 1.6}
  \begin{block}{Learning goal}
    Treat expectation as integration and deploy inequalities, change of variables, and limit theorems confidently.
  \end{block}
  \begin{enumerate}
    \item Expectation as Lebesgue integral.
    \item Linearity, monotonicity, and convergence theorems.
    \item Markov, Chebyshev, Jensen, and moment inequalities.
    \item Change of variables, tail expectations, inclusion-exclusion, and Bonferroni.
  \end{enumerate}
\end{frame}

\begin{frame}{Expected Value as Integral}
  \begin{block}{Definition}
    For a random variable $X$ on $(\Omega,\mathcal F,\Prob)$,
    \[
      \E X=\int_\Omega X\,d\Prob,
    \]
    provided the positive and negative parts do not both have infinite integral.
  \end{block}
  \begin{alertblock}{Pitfall}
    $\E X$ is not defined when $\E X^+=\E X^-=\infty$.
  \end{alertblock}
  \vikiname{Expected value as Lebesgue integral; ID: durrett_ch1_expectation_as_integral}
\end{frame}

\begin{frame}{Linearity and Monotonicity}
  \begin{block}{Theorem}
    If $X,Y$ are integrable and $a,b\in\R$, then
    \[
      \E(aX+bY)=a\E X+b\E Y.
    \]
    If $X\le Y$ a.s., then $\E X\le \E Y$ whenever the expectations are defined.
  \end{block}
  \proofblock{
  This is exactly integral linearity and monotonicity applied on the probability space. Independence is not needed for linearity; independence enters only when factoring products.
  }
  \vikiname{Linearity and monotonicity of expectation; ID: durrett_ch1_expectation_linearity}
\end{frame}

\begin{frame}{Markov and Chebyshev}
  \begin{block}{Markov}
    If $Y\ge0$ and $a>0$, then
    \[
      \Prob(Y\ge a)\le \frac{\E Y}{a}.
    \]
  \end{block}
  \begin{block}{Chebyshev}
    If $\E X=\mu$ and $\var(X)=\sigma^2$, then
    \[
      \Prob(|X-\mu|\ge a)\le \frac{\sigma^2}{a^2}.
    \]
  \end{block}
  \vikiname{Markov and Chebyshev inequalities; ID: durrett_ch1_markov_chebyshev}
\end{frame}

\begin{frame}{Markov: How to Prove}
  \proofblock{
  For $Y\ge0$,
  \[
    a\one_{\{Y\ge a\}}\le Y.
  \]
  Taking expectations gives $a\Prob(Y\ge a)\le \E Y$. Chebyshev follows by applying Markov to $Y=(X-\mu)^2$.
  }
  \vikiname{Markov and Chebyshev inequalities; ID: durrett_ch1_markov_chebyshev}
\end{frame}

\begin{frame}{Expectation Limit Theorems}
  \begin{block}{Toolkit}
    Fatou, monotone convergence, dominated convergence, and bounded convergence all transfer to expectations.
  \end{block}
  \begin{itemize}
    \item Fatou: nonnegative sequences give lower semicontinuity.
    \item MCT: $0\le X_n\uparrow X$ implies $\E X_n\uparrow \E X$.
    \item DCT: $X_n\to X$ and $|X_n|\le Y$ with $\E Y<\infty$ implies $\E X_n\to \E X$.
  \end{itemize}
  \vikiname{Limit theorems for expectation; ID: durrett_ch1_expectation_limit_theorems}
\end{frame}

\begin{frame}[shrink=8]{Change of Variables Formula}
  \begin{block}{Theorem}
    If $X$ has law $\mu$ and $f$ is measurable, then
    \[
      \E f(X)=\int f(x)\,\mu(dx),
    \]
    whenever $f\ge0$ or $f(X)$ is integrable.
  \end{block}
  \proofblock{
  Prove first for indicators $f=\one_A$, where both sides equal $\Prob(X\in A)$. Extend by linearity to simple functions, then by MCT to nonnegative $f$, and by signed decomposition to integrable $f$.
  }
  \vikiname{Change of variables formula; ID: durrett_ch1_change_of_variables}
\end{frame}

\begin{frame}{Tail Integral Formula}
  \begin{block}{Formula}
    For $X\ge0$,
    \[
      \E X=\int_0^\infty \Prob(X>t)\,dt.
    \]
  \end{block}
  \proofblock{
  Write
  \[
    X(\omega)=\int_0^\infty \one_{\{t<X(\omega)\}}\,dt.
  \]
  Then apply Tonelli to interchange the $\omega$ and $t$ integrals.
  }
  \vikiname{Tail integral and nonnegative expectations; ID: durrett_ch1_tail_sum_nonnegative}
\end{frame}

\begin{frame}{Inclusion-Exclusion}
  \begin{block}{Formula}
    For events $A_1,\ldots,A_n$,
    \[
      \Prob\left(\bigcup_i A_i\right)
      =
      \sum_i\Prob(A_i)-\sum_{i<j}\Prob(A_iA_j)+\cdots
      +(-1)^{n-1}\Prob\left(\bigcap_i A_i\right).
    \]
  \end{block}
  \proofblock{
  Use
  \[
    \one_{\cup_i A_i}=1-\prod_{i=1}^n(1-\one_{A_i}),
  \]
  expand the finite product, then take expectations.
  }
  \vikiname{Inclusion-exclusion and Bonferroni; ID: durrett_ch1_inclusion_exclusion}
\end{frame}

\begin{frame}{Exercise 1.6.1}
  Suppose $\phi$ is strictly convex. Under the assumptions of Jensen's
  inequality, show that
  \[
    \phi(EX)=E\phi(X)
  \]
  implies $X=EX$ almost surely.
  \vikiname{Equality case in strict Jensen; ID: exercise_method_strict_jensen_equality}
\end{frame}

\begin{frame}{Exercise 1.6.1: Proof Skeleton}
  \proofblock{
  Let $m=\E X$. If $\Prob(X\ne m)>0$, then $X$ is nonconstant on a set of positive probability. Strict convexity makes Jensen strict for nonconstant random variables:
  \[
    \E\phi(X)>\phi(\E X).
  \]
  This contradicts equality. Hence $\Prob(X=m)=1$.
  }
  \vikiname{Equality case in strict Jensen; ID: exercise_method_strict_jensen_equality}
\end{frame}

\begin{frame}{Exercise 1.6.2}
  Suppose $\phi:\R^n\to\R$ is convex. Show that
  \[
    E\phi(X_1,\ldots,X_n)\ge \phi(EX_1,\ldots,EX_n),
  \]
  provided $E|\phi(X_1,\ldots,X_n)|<\infty$ and $E|X_i|<\infty$ for all
  $i$.
  \vikiname{Jensen's inequality for integrals; ID: durrett_ch1_jensen_integral}
\end{frame}

\begin{frame}{Exercise 1.6.2: Proof Skeleton}
  \proofblock{
  Let $m=(\E X_1,\ldots,\E X_n)$. A convex function has a supporting hyperplane at $m$:
  \[
    \phi(x)\ge a\cdot x+b,\qquad \phi(m)=a\cdot m+b.
  \]
  Taking expectations yields
  \[
    \E\phi(X)\ge a\cdot \E X+b=\phi(m).
  \]
  }
  \vikiname{Jensen's inequality for integrals; ID: durrett_ch1_jensen_integral}
\end{frame}

\begin{frame}[shrink=6]{Exercise 1.6.3}
  Chebyshev's inequality is and is not sharp.
  \begin{enumerate}
    \item Show that if $0<b\le a$ are fixed, there is $X$ with $EX^2=b^2$
    for which $P(|X|\ge a)=b^2/a^2$.
    \item Show that if $0<EX^2<\infty$, then
    \[
      \lim_{a\to\infty}\frac{a^2P(|X|\ge a)}{EX^2}=0.
    \]
  \end{enumerate}
  \vikiname{Two-point distributions as Chebyshev extremizers; ID: exercise_method_two_point_extremizers}
\end{frame}

\begin{frame}[shrink=8]{Exercise 1.6.3: Proof Skeleton}
  \longproofblock{
  For fixed $a,b$, take
  \[
    \Prob(X=a)=\Prob(X=-a)=\frac{b^2}{2a^2},\qquad
    \Prob(X=0)=1-\frac{b^2}{a^2}.
  \]
  Then $\E X^2=b^2$ and $\Prob(|X|\ge a)=b^2/a^2$. For the limit statement set $Y=X^2$. Then
  \[
    a^2\Prob(|X|\ge a)\le \E(Y;Y\ge a^2).
  \]
  Since $Y\one_{\{Y\ge a^2\}}\to0$ and is dominated by integrable $Y$, DCT gives the limit.
  }
  \vikiname{Two-point distributions as Chebyshev extremizers; ID: exercise_method_two_point_extremizers}
\end{frame}

\begin{frame}[shrink=8]{Exercise 1.6.4}
  One-sided Chebyshev bound.
  \begin{enumerate}
    \item Let $a>b>0$, $0<p<1$, and let $X$ satisfy
    $P(X=a)=p$, $P(X=-b)=1-p$. Apply Chebyshev's inequality to
    $\phi(x)=(x+b)^2$ and conclude that if $Y$ has the same mean and
    variance as $X$, then $P(Y\ge a)\le p$, with equality when $Y=X$.
    \item If $EY=0$, $\operatorname{var}(Y)=\sigma^2$, and $a>0$, show
    \[
      P(Y\ge a)\le \frac{\sigma^2}{a^2+\sigma^2},
    \]
    and show equality is possible.
  \end{enumerate}
  \vikiname{Two-point distributions as Chebyshev extremizers; ID: exercise_method_two_point_extremizers}
\end{frame}

\begin{frame}[shrink=10]{Exercise 1.6.4: Proof Skeleton}
  \longproofblock{
  For the two-point $X$,
  \[
    \E(X+b)^2=p(a+b)^2.
  \]
  If $Y$ has the same mean and variance, then $\E(Y+b)^2$ has the same value. On $\{Y\ge a\}$,
  \[
    (Y+b)^2\ge(a+b)^2,
  \]
  so Markov gives $P(Y\ge a)\le p$. For mean zero and variance $\sigma^2$, choose the extremal two-point distribution with values $a$ and $-\sigma^2/a$. Its mass at $a$ is
  \[
    \frac{\sigma^2}{a^2+\sigma^2},
  \]
  giving both the bound and equality case.
  }
  \vikiname{Two-point distributions as Chebyshev extremizers; ID: exercise_method_two_point_extremizers}
\end{frame}

\begin{frame}[shrink=7]{Exercise 1.6.5}
  Show that the following lower bounds do not exist:
  \begin{enumerate}
    \item If $\varepsilon>0$,
    \[
      \inf\{P(|X|>\varepsilon):EX=0,\operatorname{var}(X)=1\}=0.
    \]
    \item If $y\ge1$ and $\sigma^2\in(0,\infty)$,
    \[
      \inf\{P(|X|>y):EX=1,\operatorname{var}(X)=\sigma^2\}=0.
    \]
  \end{enumerate}
  \vikiname{Two-point distributions as Chebyshev extremizers; ID: exercise_method_two_point_extremizers}
\end{frame}

\begin{frame}[shrink=8]{Exercise 1.6.5: Proof Skeleton}
  \longproofblock{
  Fix $\delta>0$ and put most mass at the mean. For part (i), take
  \[
    \Prob(X=0)=1-\delta,\qquad
    \Prob(X=\delta^{-1/2})=\Prob(X=-\delta^{-1/2})=\delta/2.
  \]
  Then $\E X=0$ and $\var(X)=1$, while the tail event has probability at most $\delta$. For part (ii), use $X=1+\sigma Z$, where $Z$ is the variable from part (i). Then $\E X=1$, $\var(X)=\sigma^2$, and the event $\{|X|>y\}$ is contained in the rare nonzero part for small $\delta$.
  }
  \vikiname{Two-point distributions as Chebyshev extremizers; ID: exercise_method_two_point_extremizers}
\end{frame}

\begin{frame}{Exercise 1.6.6}
  Let $Y\ge0$ with $EY^2<\infty$. Apply Cauchy--Schwarz to
  $Y\mathbf{1}_{\{Y>0\}}$ and conclude
  \[
    P(Y>0)\ge \frac{(EY)^2}{EY^2}.
  \]
  \vikiname{Cauchy-Schwarz lower bound for positive probability; ID: exercise_method_paley_zygmund_basic}
\end{frame}

\begin{frame}{Exercise 1.6.6: Proof Skeleton}
  \proofblock{
  Since $Y=Y\one_{\{Y>0\}}$,
  \[
    \E Y=\E(Y\one_{\{Y>0\}}).
  \]
  Cauchy--Schwarz gives
  \[
    (\E Y)^2\le \E(Y^2)\E(\one_{\{Y>0\}}^2)=\E Y^2\,\Prob(Y>0).
  \]
  Divide by $\E Y^2$ unless it is zero, in which case the claim is trivial.
  }
  \vikiname{Cauchy-Schwarz lower bound for positive probability; ID: exercise_method_paley_zygmund_basic}
\end{frame}

\begin{frame}[shrink=8]{Exercise 1.6.7}
  Let $\Omega=(0,1)$ with Lebesgue measure, let $\alpha\in(1,2)$, and set
  \[
    X_n=n^\alpha\mathbf{1}_{(1/(n+1),1/n)}.
  \]
  Show $X_n\to0$ almost surely and Theorem 1.6.8 applies with
  $h(x)=x$ and $g(x)=|x|^{2/\alpha}$, but the $X_n$ are not dominated by
  an integrable random variable.
  \vikiname{Limit theorems for expectation; ID: durrett_ch1_expectation_limit_theorems}
\end{frame}

\begin{frame}[shrink=10]{Exercise 1.6.7: Proof Skeleton}
  \longproofblock{
  The intervals $(1/(n+1),1/n)$ are disjoint, so each $\omega$ belongs to at most one of them; hence $X_n(\omega)\to0$. Also
  \[
    \E g(X_n)=n^2\left(\frac1n-\frac1{n+1}\right)=\frac{n}{n+1}\le1,
  \]
  and $|h(x)|/g(x)=|x|^{1-2/\alpha}\to0$. If $|X_n|\le Y$ with $\E Y<\infty$, then on $(1/(n+1),1/n)$, $Y\ge n^\alpha$, so
  \[
    \E Y\ge \sum_{n=1}^N n^\alpha\left(\frac1n-\frac1{n+1}\right),
  \]
  whose terms are comparable to $n^{\alpha-2}$; the series diverges since $\alpha\in(1,2)$.
  }
  \vikiname{Limit theorems for expectation; ID: durrett_ch1_expectation_limit_theorems}
\end{frame}

\begin{frame}{Exercise 1.6.8}
  Suppose a probability measure $\mu$ has
  \[
    \mu(A)=\int_A f(x)\,dx
  \]
  for all Borel $A$. Show that for any $g$ with $g\ge0$ or
  $\int |g|\,d\mu<\infty$,
  \[
    \int g(x)\,\mu(dx)=\int g(x)f(x)\,dx.
  \]
  \vikiname{Change of variables formula; ID: durrett_ch1_change_of_variables}
\end{frame}

\begin{frame}{Exercise 1.6.8: Proof Skeleton}
  \proofblock{
  First prove it for $g=\one_A$:
  \[
    \int \one_A\,d\mu=\mu(A)=\int_A f(x)\,dx.
  \]
  Extend by linearity to nonnegative simple functions. For $g\ge0$, take simple $g_n\uparrow g$ and apply MCT on both sides. For signed integrable $g$, apply the nonnegative result to $g^+$ and $g^-$ and subtract.
  }
  \vikiname{Change of variables formula; ID: durrett_ch1_change_of_variables}
\end{frame}

\begin{frame}{Exercise 1.6.9}
  Let $A_1,\ldots,A_n$ be events and $A=\bigcup_{i=1}^n A_i$. Prove the
  inclusion--exclusion formula.
  \vikiname{Indicator expansion for inclusion-exclusion and Bonferroni; ID: exercise_method_inclusion_exclusion_indicators}
\end{frame}

\begin{frame}[shrink=8]{Exercise 1.6.9: Proof Skeleton}
  \longproofblock{
  Pointwise,
  \[
    \one_A=1-\prod_{i=1}^n(1-\one_{A_i}).
  \]
  Expanding the product gives
  \[
    \one_A=\sum_i\one_{A_i}-\sum_{i<j}\one_{A_i\cap A_j}
    +\sum_{i<j<k}\one_{A_i\cap A_j\cap A_k}-\cdots.
  \]
  Taking expectations yields the finite inclusion--exclusion identity.
  }
  \vikiname{Indicator expansion for inclusion-exclusion and Bonferroni; ID: exercise_method_inclusion_exclusion_indicators}
\end{frame}

\begin{frame}[shrink=8]{Exercise 1.6.10}
  Prove the Bonferroni inequalities for
  $A=\bigcup_{i=1}^n A_i$: stopping inclusion--exclusion after an even
  number of sums gives a lower bound, and stopping after an odd number of
  sums gives an upper bound.
  \vikiname{Indicator expansion for inclusion-exclusion and Bonferroni; ID: exercise_method_inclusion_exclusion_indicators}
\end{frame}

\begin{frame}[shrink=8]{Exercise 1.6.10: Proof Skeleton}
  \longproofblock{
  For each outcome let
  \[
    N=\sum_i\one_{A_i}.
  \]
  Then $\one_A=\one_{\{N\ge1\}}$. The inclusion--exclusion partial sums become
  \[
    S_r(N)=\sum_{k=1}^r(-1)^{k-1}\binom{N}{k}.
  \]
  For $N=0$ all terms vanish. For $N\ge1$, alternating binomial partial sums lie alternately above and below $1$. Thus odd truncations dominate $\one_A$, even truncations are below it. Take expectations.
  }
  \vikiname{Indicator expansion for inclusion-exclusion and Bonferroni; ID: exercise_method_inclusion_exclusion_indicators}
\end{frame}

\begin{frame}{Exercise 1.6.11}
  If $E|X|^k<\infty$, show that for $0<j<k$,
  \[
    E|X|^j<\infty,
    \qquad
    E|X|^j\le (E|X|^k)^{j/k}.
  \]
  \vikiname{Jensen's inequality for integrals; ID: durrett_ch1_jensen_integral}
\end{frame}

\begin{frame}{Exercise 1.6.11: Proof Skeleton}
  \proofblock{
  Let $Y=|X|^k$ and $r=j/k\in(0,1)$. The map $t\mapsto t^r$ is concave, so Jensen gives
  \[
    \E Y^r\le (\E Y)^r.
  \]
  Since $Y^r=|X|^j$,
  \[
    \E|X|^j\le(\E|X|^k)^{j/k}<\infty.
  \]
  }
  \vikiname{Jensen's inequality for integrals; ID: durrett_ch1_jensen_integral}
\end{frame}

\begin{frame}[shrink=6]{Exercise 1.6.12}
  Apply Jensen's inequality with $\phi(x)=e^x$ and
  $P(X=\log y_m)=p(m)$ to prove
  \[
    \sum_{m=1}^n p(m)y_m
    \ge
    \prod_{m=1}^n y_m^{p(m)}
  \]
  when $\sum_m p(m)=1$ and $p(m),y_m>0$.
  \vikiname{Jensen's inequality for integrals; ID: durrett_ch1_jensen_integral}
\end{frame}

\begin{frame}{Exercise 1.6.12: Proof Skeleton}
  \proofblock{
  Let $\Prob(X=\log y_m)=p(m)$. Jensen for $\phi(x)=e^x$ gives
  \[
    \E e^X\ge e^{\E X}.
  \]
  Now
  \[
    \E e^X=\sum_{m=1}^n p(m)y_m,
    \qquad
    e^{\E X}=\exp\left(\sum_m p(m)\log y_m\right)=\prod_m y_m^{p(m)}.
  \]
  }
  \vikiname{Jensen's inequality for integrals; ID: durrett_ch1_jensen_integral}
\end{frame}

\begin{frame}{Exercise 1.6.13}
  If $EX_1^-<\infty$ and $X_n\uparrow X$, show that
  \[
    EX_n\uparrow EX.
  \]
  \vikiname{Limit theorems for expectation; ID: durrett_ch1_expectation_limit_theorems}
\end{frame}

\begin{frame}{Exercise 1.6.13: Proof Skeleton}
  \proofblock{
  Let $Y_n=X_n-X_1$. Then $Y_n\ge0$ and $Y_n\uparrow X-X_1$. By MCT,
  \[
    \E Y_n\uparrow \E(X-X_1).
  \]
  Since $\E X_1^-<\infty$, adding $\E X_1$ is legitimate in the extended sense:
  \[
    \E X_n=\E X_1+\E Y_n\uparrow \E X.
  \]
  }
  \vikiname{Limit theorems for expectation; ID: durrett_ch1_expectation_limit_theorems}
\end{frame}

\begin{frame}{Exercise 1.6.14}
  Let $X\ge0$, but do not assume $E(1/X)<\infty$. Show
  \[
    \lim_{y\to\infty} yE(1/X;X>y)=0,
    \qquad
    \lim_{y\downarrow0} yE(1/X;X>y)=0.
  \]
  \vikiname{Tail truncation for expectation limits; ID: exercise_method_expectation_tail_truncation}
\end{frame}

\begin{frame}[shrink=8]{Exercise 1.6.14: Proof Skeleton}
  \longproofblock{
  As $y\to\infty$,
  \[
    0\le y\E(1/X;X>y)=\E(y/X;X>y)\le \Prob(X>y)\to0.
  \]
  As $y\downarrow0$, fix $\eta>0$ and choose $\delta>0$ such that $\Prob(0<X\le\delta)<\eta$. For $0<y<\delta$,
  \[
    y\E(1/X;X>y)
    \le \Prob(y<X\le\delta)+y/\delta
    \le \eta+y/\delta.
  \]
  Let $y\downarrow0$, then $\eta\downarrow0$.
  }
  \vikiname{Tail truncation for expectation limits; ID: exercise_method_expectation_tail_truncation}
\end{frame}

\begin{frame}{Exercise 1.6.15}
  If $X_n\ge0$, show that
  \[
    E\left(\sum_{n=0}^{\infty}X_n\right)=\sum_{n=0}^{\infty}EX_n.
  \]
  \vikiname{Limit theorems for expectation; ID: durrett_ch1_expectation_limit_theorems}
\end{frame}

\begin{frame}{Exercise 1.6.15: Proof Skeleton}
  \proofblock{
  Let
  \[
    S_N=\sum_{n=0}^N X_n.
  \]
  Then $S_N\uparrow \sum_{n=0}^\infty X_n$. MCT gives
  \[
    \E\sum_{n=0}^\infty X_n=\lim_{N\to\infty}\E S_N.
  \]
  Finite linearity gives $\E S_N=\sum_{n=0}^N\E X_n$; let $N\to\infty$.
  }
  \vikiname{Limit theorems for expectation; ID: durrett_ch1_expectation_limit_theorems}
\end{frame}

\begin{frame}{Exercise 1.6.16}
  If $X$ is integrable and $A_n$ are disjoint sets with union $A$, show
  \[
    \sum_{n=0}^{\infty}E(X;A_n)=E(X;A),
  \]
  i.e. the sum converges absolutely and has the displayed value.
  \vikiname{Linearity and monotonicity of expectation; ID: durrett_ch1_expectation_linearity}
\end{frame}

\begin{frame}[shrink=8]{Exercise 1.6.16: Proof Skeleton}
  \longproofblock{
  Since $X$ is integrable,
  \[
    \sum_n \E(|X|;A_n)=\E(|X|;A)\le \E|X|<\infty,
  \]
  so $\sum_n\E(X;A_n)$ converges absolutely. For finite unions,
  \[
    \sum_{n=0}^N\E(X;A_n)=\E\left(X;\bigcup_{n=0}^N A_n\right).
  \]
  As $N\to\infty$, $X\one_{\cup_{n=0}^N A_n}\to X\one_A$ and is dominated by $|X|$. DCT gives the limit.
  }
  \vikiname{Linearity and monotonicity of expectation; ID: durrett_ch1_expectation_linearity}
\end{frame}

\end{document}
